
%% ====================== ปก ======================
\begin{frame}[plain]
  \maketitle
\end{frame}

%% ====================== ทบทวน ======================
\begin{frame}{ทบทวนสัปดาห์ที่ 2}
  \begin{itemize}
    \item \textbf{ประเภทประพจน์}:
          \begin{itemize}
            \item \textbf{สัจนิรันดร์} (Tautology) — จริงทุกกรณี เช่น $p \lor \neg p$
            \item \textbf{ข้อขัดแย้ง} (Contradiction) — เท็จทุกกรณี เช่น $p \land \neg p$
            \item \textbf{ประพจน์กลาง} (Contingency) — จริงบ้างเท็จบ้าง เช่น $p \rightarrow q$
          \end{itemize}
    \item \textbf{ความสมมูลทางตรรกศาสตร์ ($P \equiv Q$)}: มีค่าความจริงเหมือนกันทุกกรณี
    \item \textbf{กฎพีชคณิตของประพจน์}: กฎเอกลักษณ์ กฎการสลับที่ กฎการแจกแจง และกฎเดอมอร์แกน
    \item \textbf{ประพจน์เงื่อนไข $p \rightarrow q$}: สมมูลกับ $\neg q \rightarrow \neg p$
  \end{itemize}
\end{frame}

%% ====================== วัตถุประสงค์ ======================
\begin{frame}{วัตถุประสงค์การเรียนรู้ (สัปดาห์ที่ 3)}
  เมื่อเรียนจบสัปดาห์นี้ นักศึกษาจะสามารถ
  \begin{enumerate}
    \item เข้าใจและใช้ตัวเชื่อมพิเศษ \textbf{XOR ($\oplus$), NAND ($\uparrow$), NOR ($\downarrow$)} ได้
    \item หาค่าความจริงและเขียนนิเสธของประพจน์ที่มี \textbf{ตัวบ่งปริมาณ} ($\forall, \exists$) ได้
    \item ประยุกต์ใช้ \textbf{กฎการอนุมาน (Rules of Inference)} เพื่อหาผลสรุปที่สมเหตุสมผลได้
    \item ตรวจสอบความ \textbf{สมเหตุสมผลของการอ้างเหตุผล} ด้วยตารางค่าความจริงได้
  \end{enumerate}
\end{frame}

%% ====================== หัวข้อ ======================
\begin{frame}{หัวข้อที่จะศึกษาในสัปดาห์นี้}
  \tableofcontents
\end{frame}


%% =====================================================================
\section{ตัวเชื่อมเพิ่มเติม: XOR, NAND, NOR}
%% =====================================================================

\begin{frame}{ประพจน์เลือกเฉพาะ (XOR, $\oplus$)}
  \begin{columns}[T]
    \begin{column}{0.5\textwidth}
      $p \oplus q$ จริง \textbf{ก็ต่อเมื่อค่าต่างกัน} \\
      \medskip
      \begin{center}
      \begin{tabular}{|c|c|c|}
        \hline \textbf{$p$} & \textbf{$q$} & \textbf{$p\oplus q$} \\ \hline
        \TT & \TT & \FF \\ \hline
        \TT & \FF & \TT \\ \hline
        \FF & \TT & \TT \\ \hline
        \FF & \FF & \FF \\ \hline
      \end{tabular}
      \end{center}
    \end{column}
    \begin{column}{0.5\textwidth}
      \centering
      \begin{tikzpicture}[scale=0.8]
        \node[xor gate US, draw, logic gate inputs=nn, fill=lightnavy, minimum height=1.2cm] (X) at (0,0) {};
        \draw (X.input 1) -- ++(-0.8,0) node[left]{$p$};
        \draw (X.input 2) -- ++(-0.8,0) node[left]{$q$};
        \draw (X.output) -- ++(0.8,0) node[right]{$p\oplus q$};
      \end{tikzpicture}
      \par\medskip
      \footnotesize $p\oplus q \equiv (p\land\neg q)\lor(\neg p\land q)$
    \end{column}
  \end{columns}
  \vspace{0.4em}
  \begin{exampleblock}{การประยุกต์ในวิทยาการคอมพิวเตอร์}
    Half-Adder (บิตผลรวม), การเข้ารหัสลับ (One-Time Pad: $M=C\oplus K$),
    Parity Check, และการสลับค่าตัวแปรโดยไม่ใช้ตัวแปรชั่วคราว
  \end{exampleblock}
\end{frame}

\begin{frame}{แนนด์และนอร์ (NAND $\uparrow$, NOR $\downarrow$)}
  \begin{block}{นิยาม}
    $p \uparrow q \equiv \neg(p \land q)$ \quad (NAND, Sheffer stroke) \\
    $p \downarrow q \equiv \neg(p \lor q)$ \quad (NOR, Peirce arrow)
  \end{block}
  \begin{center}\small
  \begin{tabular}{|c|c|c|c|}
    \hline
    \textbf{$p$} & \textbf{$q$} & \textbf{$p\uparrow q$ (NAND)} & \textbf{$p\downarrow q$ (NOR)} \\ \hline
    \TT & \TT & \FF & \FF \\ \hline
    \TT & \FF & \TT & \FF \\ \hline
    \FF & \TT & \TT & \FF \\ \hline
    \FF & \FF & \TT & \TT \\ \hline
  \end{tabular}
  \end{center}
  \vspace{0.3em}
  \begin{exampleblock}{ตัวเชื่อมสากล (Universal Gate)}
    ใช้ NAND อย่างเดียว (หรือ NOR อย่างเดียว) สร้าง NOT, AND, OR ได้ทั้งหมด \\
    เช่น $\neg p \equiv p \uparrow p$ \quad ; \quad $p \land q \equiv (p\uparrow q)\uparrow(p\uparrow q)$
  \end{exampleblock}
\end{frame}

\begin{frame}{ตัวเชื่อมสากล (Universal Gate) -- สร้างทุกตัวเชื่อมจาก NAND}
  \begin{block}{ทฤษฎีบท}
    NAND ($\uparrow$) อย่างเดียว (หรือ NOR อย่างเดียว) สามารถสร้าง \textbf{ทุกตัวเชื่อม} ได้
    จึงเรียกว่า \textbf{ตัวเชื่อมสากล (Universal Gate)} -- เป็นพื้นฐานของวงจรดิจิทัล
  \end{block}
  \vspace{0.3em}
  \begin{center}
  \begin{tabular}{|l|l|}
    \hline
    \textbf{ตัวเชื่อม} & \textbf{เขียนด้วย NAND อย่างเดียว} \\ \hline
    $\neg p$ & $p \uparrow p$ \\ \hline
    $p \land q$ & $(p \uparrow q) \uparrow (p \uparrow q)$ \\ \hline
    $p \lor q$ & $(p \uparrow p) \uparrow (q \uparrow q)$ \\ \hline
    $p \rightarrow q$ & $p \uparrow (q \uparrow q)$ \\ \hline
  \end{tabular}
  \end{center}
  \vspace{0.4em}
  \begin{exampleblock}{ตัวอย่างการพิสูจน์: $p \lor q \equiv (p \uparrow p) \uparrow (q \uparrow q)$}
    \footnotesize
    $p \lor q \equiv \neg(\neg p \land \neg q)$ \quad (กฎเดอมอร์แกน) \\
    $\equiv \neg(p \uparrow p) \land (q \uparrow q)$ \quad (เพราะ $\neg p \equiv p \uparrow p$) \\
    $\equiv (p \uparrow p) \uparrow (q \uparrow q)$ \quad (นิยาม NAND)
  \end{exampleblock}
\end{frame}


%% =====================================================================
\section{ประพจน์เปิดและตัวบ่งปริมาณ}
%% =====================================================================

\begin{frame}{ตัวบ่งปริมาณ (Quantifiers)}
  ประพจน์เปิด เช่น $P(x): ``x > 5''$ ระบุค่าความจริงไม่ได้จนกว่าจะแทนค่า \\
  ทำให้ ``ปิด'' ได้ด้วย \textbf{ตัวบ่งปริมาณ} 2 ชนิด
  \begin{block}{}
    \begin{itemize}
      \item $\forall x\,P(x)$ -- ตัวบ่งปริมาณ \textbf{สำหรับทั้งหมด} (Universal, ``สำหรับทุก'') \\
            หมายถึง $P(x)$ จริงสำหรับ $x$ \emph{ทุกตัว} ในเอกภพสัมพัทธ์
      \item $\exists x\,P(x)$ -- ตัวบ่งปริมาณ \textbf{สำหรับบางตัว} (Existential, ``มีบางตัว'') \\
            หมายถึง มี $x$ \emph{อย่างน้อยหนึ่งตัว} ที่ทำให้ $P(x)$ จริง
    \end{itemize}
  \end{block}
  \begin{alertblock}{เทคนิค}
    แสดงว่า $\forall x\,P(x)$ เท็จ $\Rightarrow$ ยก \textbf{ตัวอย่างค้าน} เพียงตัวเดียว \\
    แสดงว่า $\exists x\,P(x)$ จริง $\Rightarrow$ ยกตัวอย่างที่ใช้ได้เพียงตัวเดียว
  \end{alertblock}
\end{frame}

\begin{frame}{ตัวอย่าง: การตีความตัวบ่งปริมาณ}
  กำหนดเอกภพสัมพัทธ์เป็นจำนวนจริง $\mathbb{R}$
  \begin{enumerate}
    \item $\forall x\,(x^2 \geq 0)$ \hfill \textcolor{Tgreen}{\textbf{จริง}} -- กำลังสองของจำนวนจริงใด ๆ ไม่ติดลบ
    \item $\exists x\,(x^2 = 4)$ \hfill \textcolor{Tgreen}{\textbf{จริง}} -- มี $x=2$ (หรือ $x=-2$)
    \item $\forall x\,(x^2 = 4)$ \hfill \textcolor{Fred}{\textbf{เท็จ}} -- ตัวอย่างค้าน $x=0$ ได้ $x^2=0\neq4$
  \end{enumerate}
  \vspace{0.8em}
  \begin{block}{นิเสธของตัวบ่งปริมาณ}
    \[
      \neg(\forall x\,P(x)) \equiv \exists x\,\neg P(x), \qquad
      \neg(\exists x\,P(x)) \equiv \forall x\,\neg P(x)
    \]
    เช่น นิเสธของ ``นักศึกษาทุกคนสอบผ่าน'' คือ ``มีนักศึกษาอย่างน้อยหนึ่งคนสอบไม่ผ่าน''
  \end{block}
\end{frame}

\begin{frame}{การยกตัวอย่างค้าน (Counterexample)}
  \begin{alertblock}{หลักการสำคัญ}
    \begin{itemize}
      \item แสดงว่า $\forall x\,P(x)$ \textbf{เท็จ} $\Rightarrow$ หา \textbf{ตัวอย่างค้าน} เพียง \textbf{ตัวเดียว} ที่ทำให้ $P(x)$ เป็นเท็จ
      \item แสดงว่า $\exists x\,P(x)$ \textbf{จริง} $\Rightarrow$ หาตัวอย่างเพียง \textbf{ตัวเดียว} ที่ทำให้ $P(x)$ เป็นจริง
    \end{itemize}
  \end{alertblock}
  \vspace{0.4em}
  \begin{exampleblock}{ตัวอย่าง: หาตัวอย่างค้าน}
    เอกภพ $U = \{1, 2, 3, 4, 5\}$, ประพจน์ $\forall x\,(x > 3)$
    \begin{itemize}
      \item ตรวจ: $1>3$ = \FF, $2>3$ = \FF, \ldots\ $\Rightarrow$ เท็จ
      \item \textbf{ตัวอย่างค้าน: $x = 1$} (เพราะ $1 \not> 3$)
    \end{itemize}
  \end{exampleblock}
  \begin{exampleblock}{ตัวอย่าง: หาตัวอย่างยืนยัน}
    เอกภพ $U = \{1, 2, 3, 4, 5\}$, ประพจน์ $\exists x\,(x \text{ เป็นจำนวนคู่})$ \\
    \textbf{ตัวอย่างยืนยัน: $x = 2$} (เพราะ 2 เป็นจำนวนคู่)
  \end{exampleblock}
\end{frame}


%% =====================================================================
\section{การอนุมาน (Rules of Inference)}
%% =====================================================================

\begin{frame}{การอนุมานคืออะไร}
  \begin{block}{นิยาม}
    \textbf{การอนุมาน (Inference Rules)} คือ กฎสำหรับสรุปผลที่สมเหตุสมผลจากเหตุที่กำหนด
    เป็นเครื่องมือพื้นฐานของ \textbf{การให้เหตุผลแบบนิรนัย (Deductive Reasoning)}
  \end{block}
  \vspace{0.4em}
  \begin{itemize}
    \item เขียนในรูป \textbf{ข้อตั้งต้น (Premises)} และ \textbf{ผลสรุป (Conclusion)}
    \item รูปแบบที่ ``สมเหตุสมผล'' = ประพจน์ (ข้อตั้งต้นทั้งหมด $\land$) $\rightarrow$ ผลสรุป เป็น \textbf{สัจนิรันดร์}
  \end{itemize}
  \vspace{0.4em}
  รูปแบบที่จะศึกษา: Modus Ponens, Modus Tollens, Hypothetical Syllogism, Disjunctive Syllogism
\end{frame}

\begin{frame}{กฎการยืนยัน (Modus Ponens)}
  \begin{columns}[T]
    \begin{column}{0.42\textwidth}
      \begin{block}{รูปแบบ}
        \begin{tabular}{l l}
          เหตุ 1: & $p \rightarrow q$ \\
          เหตุ 2: & $p$ \\ \hline
          ผล: & $\therefore q$ \\
        \end{tabular}
      \end{block}
      ``ถ้า $p$ แล้ว $q$ และ $p$ จริง \\ ดังนั้น $q$ จริง''
    \end{column}
    \begin{column}{0.56\textwidth}
      พิสูจน์ว่า $((p\rightarrow q)\land p)\rightarrow q$ เป็นสัจนิรันดร์
      \begin{center}\scriptsize
      \begin{tabular}{|c|c|c|c|c|}
        \hline
        $p$ & $q$ & $p\rightarrow q$ & $(p\rightarrow q)\land p$ & ผล$\rightarrow q$ \\ \hline
        \TT & \TT & \TT & \TT & \TT \\ \hline
        \TT & \FF & \FF & \FF & \TT \\ \hline
        \FF & \TT & \TT & \FF & \TT \\ \hline
        \FF & \FF & \TT & \FF & \TT \\ \hline
      \end{tabular}
      \end{center}
    \end{column}
  \end{columns}
  \vspace{0.4em}
  \begin{exampleblock}{ตัวอย่าง}
    เหตุ: ถ้าฝนตกแล้วพื้นถนนเปียก ; ฝนตก \quad
    $\Rightarrow$ \textbf{$\therefore$ พื้นถนนเปียก}
  \end{exampleblock}
\end{frame}

\begin{frame}{Modus Ponens: ตัวอย่างเพิ่มเติมและการเขียนโปรแกรม}
  \begin{exampleblock}{ตัวอย่าง}
    เหตุ 1: ถ้าคนรักษาสัญญา แล้วคนนั้นจะได้รับความเชื่อถือ \\
    เหตุ 2: นาย ก รักษาสัญญา \\
    \textbf{$\therefore$ นาย ก ได้รับความเชื่อถือ}
  \end{exampleblock}
  \vspace{0.6em}
  \begin{block}{เชื่อมโยงกับการเขียนโปรแกรม}
    คำสั่ง \texttt{if (condition) \{ execute\_action(); \}} \\
    ทำงานบนหลักเดียวกับ Modus Ponens: ถ้าเงื่อนไขจริง แล้วกระทำการที่กำหนด
  \end{block}
\end{frame}

\begin{frame}{กฎการปฏิเสธ (Modus Tollens)}
  \begin{columns}[T]
    \begin{column}{0.42\textwidth}
      \begin{block}{รูปแบบ}
        \begin{tabular}{l l}
          เหตุ 1: & $p \rightarrow q$ \\
          เหตุ 2: & $\neg q$ \\ \hline
          ผล: & $\therefore \neg p$ \\
        \end{tabular}
      \end{block}
      ``ถ้า $p$ แล้ว $q$ และ $q$ เท็จ \\ ดังนั้น $p$ เท็จ''
    \end{column}
    \begin{column}{0.56\textwidth}
      พิสูจน์ว่า $((p\rightarrow q)\land\neg q)\rightarrow\neg p$ เป็นสัจนิรันดร์
      \begin{center}\scriptsize
      \begin{tabular}{|c|c|c|c|c|}
        \hline
        $p$ & $q$ & $\neg q$ & $(p\rightarrow q)\land\neg q$ & ผล$\rightarrow\neg p$ \\ \hline
        \TT & \TT & \FF & \FF & \TT \\ \hline
        \TT & \FF & \TT & \FF & \TT \\ \hline
        \FF & \TT & \FF & \FF & \TT \\ \hline
        \FF & \FF & \TT & \TT & \TT \\ \hline
      \end{tabular}
      \end{center}
    \end{column}
  \end{columns}
  \vspace{0.4em}
  \begin{exampleblock}{ตัวอย่าง}
    เหตุ: ถ้าฝนตกแล้วพื้นถนนเปียก ; พื้นถนนไม่เปียก \quad
    $\Rightarrow$ \textbf{$\therefore$ ฝนไม่ตก}
  \end{exampleblock}
\end{frame}

\begin{frame}{Modus Tollens: ตัวอย่างเพิ่มเติม}
  \begin{exampleblock}{ตัวอย่าง}
    เหตุ 1: ถ้าสัตว์เป็นสุนัข แล้วสัตว์นั้นเป็นสัตว์เลี้ยงลูกด้วยนม \\
    เหตุ 2: สัตว์ตัวนี้ไม่เป็นสัตว์เลี้ยงลูกด้วยนม \\
    \textbf{$\therefore$ สัตว์ตัวนี้ไม่ใช่สุนัข}
  \end{exampleblock}
  \vspace{0.6em}
  \begin{alertblock}{ข้อควรระวังเชิงวิทยาศาสตร์}
    การสรุปว่า $p$ เท็จนั้น \emph{สมเหตุสมผลทางตรรกะ} \\
    แต่ในโลกจริงต้องระวังว่าอาจมีสาเหตุอื่นที่ทำให้ $q$ ไม่เกิด
    (เช่น พื้นเปียกเพราะรดน้ำต้นไม้)
  \end{alertblock}
\end{frame}

\begin{frame}{Hypothetical Syllogism: พิสูจน์ด้วยตารางค่าความจริง}
  \begin{block}{รูปแบบ}
    \begin{tabular}{l l}
      เหตุ 1: & $p \rightarrow q$ \\
      เหตุ 2: & $q \rightarrow r$ \\ \hline
      ผล: & $\therefore p \rightarrow r$ \\
    \end{tabular}
  \end{block}
  พิสูจน์ว่า $[(p\rightarrow q) \land (q\rightarrow r)] \rightarrow (p\rightarrow r)$ เป็นสัจนิรันดร์
  \begin{center}\footnotesize
  \begin{tabular}{|c|c|c|c|c|c|c|}
    \hline
    $p$ & $q$ & $r$ & $p\rightarrow q$ & $q\rightarrow r$ & $(p\rightarrow q)\land(q\rightarrow r)$ & $p\rightarrow r$ \\ \hline
    \TT & \TT & \TT & \TT & \TT & \TT & \TT \\ \hline
    \TT & \TT & \FF & \TT & \FF & \FF & \FF \\ \hline
    \TT & \FF & \TT & \FF & \TT & \FF & \TT \\ \hline
    \TT & \FF & \FF & \FF & \TT & \FF & \FF \\ \hline
    \FF & \TT & \TT & \TT & \TT & \TT & \TT \\ \hline
    \FF & \TT & \FF & \TT & \FF & \FF & \TT \\ \hline
    \FF & \FF & \TT & \TT & \TT & \TT & \TT \\ \hline
    \FF & \FF & \FF & \TT & \TT & \TT & \TT \\ \hline
  \end{tabular}
  \end{center}
  \vspace{0.3em}
  \begin{exampleblock}{ผลลัพธ์}
    คอลัมน์สุดท้ายเป็น \TT\ ทุกแถว $\Rightarrow$ \textbf{สัจนิรันดร์} $\Rightarrow$ กฎ \textbf{HS สมเหตุสมผล}
  \end{exampleblock}
\end{frame}

\begin{frame}{Hypothetical Syllogism: ตัวอย่าง}
  หลักการ \textbf{การถ่ายโอนเงื่อนไข (Transitivity)}: ถ้า $p$ นำสู่ $q$ และ $q$ นำสู่ $r$ แล้ว $p$ นำสู่ $r$
  \vspace{0.4em}
  \begin{exampleblock}{ตัวอย่าง}
    เหตุ 1: ถ้าฝนตก แล้วถนนจะเปียก \\
    เหตุ 2: ถ้าถนนเปียก แล้วรถจะลื่น \\
    \textbf{$\therefore$ ถ้าฝนตก แล้วรถจะลื่น}
  \end{exampleblock}
  \begin{exampleblock}{ตัวอย่างในการเขียนโปรแกรม}
    input ถูกต้อง $\rightarrow$ process ทำงาน $\rightarrow$ output ถูกสร้าง \\
    ถ้า input ถูกต้อง $\rightarrow$ output ถูกสร้าง (HS)
  \end{exampleblock}
\end{frame}

\begin{frame}{Disjunctive Syllogism: พิสูจน์ด้วยตารางค่าความจริง}
  \begin{block}{รูปแบบ}
    \begin{tabular}{l l}
      เหตุ 1: & $p \lor q$ \\
      เหตุ 2: & $\neg p$ \\ \hline
      ผล: & $\therefore q$ \\
    \end{tabular}
  \end{block}
  พิสูจน์ว่า $[(p \lor q) \land \neg p] \rightarrow q$ เป็นสัจนิรันดร์
  \begin{center}\footnotesize
  \begin{tabular}{|c|c|c|c|c|}
    \hline
    $p$ & $q$ & $p \lor q$ & $\neg p$ & $(p \lor q) \land \neg p$ \\ \hline
    \TT & \TT & \TT & \FF & \FF \\ \hline
    \TT & \FF & \TT & \FF & \FF \\ \hline
    \FF & \TT & \TT & \TT & \TT \\ \hline
    \FF & \FF & \FF & \TT & \FF \\ \hline
  \end{tabular}
  \end{center}
  \begin{center}\footnotesize
  \begin{tabular}{|c|c|}
    \hline
    $(p \lor q) \land \neg p$ & $[(p \lor q) \land \neg p] \rightarrow q$ \\ \hline
    \FF & \TT \\ \hline
    \FF & \TT \\ \hline
    \TT & \TT \\ \hline
    \FF & \TT \\ \hline
  \end{tabular}
  \end{center}
  \vspace{0.3em}
  \begin{exampleblock}{ผลลัพธ์}
    คอลัมน์สุดท้ายเป็น \TT\ ทุกแถว $\Rightarrow$ \textbf{สัจนิรันดร์} $\Rightarrow$ กฎ \textbf{DS สมเหตุสมผล}
  \end{exampleblock}
\end{frame}

\begin{frame}{Disjunctive Syllogism: ตัวอย่าง}
  หลักการ: รู้ว่าอย่างน้อยหนึ่งตัวจริง และตัวหนึ่งเท็จ $\Rightarrow$ อีกตัวต้องจริง
  \vspace{0.4em}
  \begin{exampleblock}{ตัวอย่าง}
    เหตุ 1: นาย ก อยู่บ้านหรืออยู่ที่ทำงาน \\
    เหตุ 2: นาย ก ไม่ได้อยู่ที่ทำงาน \\
    \textbf{$\therefore$ นาย ก อยู่บ้าน}
  \end{exampleblock}
  \begin{exampleblock}{ตัวอย่างในการเขียนโปรแกรม}
    error เป็น E1 หรือ E2 ; ไม่ใช่ E1 \\
    $\therefore$ error เป็น E2 (DS)
  \end{exampleblock}
\end{frame}

\begin{frame}{ข้อผิดพลาดที่พบบ่อย: Fallacies}
  \begin{alertblock}{Affirming the Consequent (ยอมรับผลที่ตามมา)}
    \begin{tabular}{l l}
      เหตุ 1: & $p \rightarrow q$ \\
      เหตุ 2: & $q$ \\ \hline
      ผลที่ผิด: & $\therefore p$ \quad (ผิด!)
    \end{tabular} \\
    ``ถ้าฝนตกแล้วดินเปียก ; ดินเปียก $\Rightarrow$ ฝนตก'' (อาจจะรดน้ำต้นไม้!)
  \end{alertblock}
  \vspace{0.4em}
  \begin{alertblock}{Denying the Antecedent (ปฏิเสธเหตุ)}
    \begin{tabular}{l l}
      เหตุ 1: & $p \rightarrow q$ \\
      เหตุ 2: & $\neg p$ \\ \hline
      ผลที่ผิด: & $\therefore \neg q$ \quad (ผิด!)
    \end{tabular} \\
    ``ถ้าเป็นนกแล้วบินได้ ; ไม่ใช่นก $\Rightarrow$ บินไม่ได้'' (เครื่องบินบินได้!)
  \end{alertblock}
  \vspace{0.4em}
  \begin{block}{ข้อสังเกต}
    ทั้งสองแบบ \textbf{ไม่ใช่สัจนิรันดร์} เมื่อตรวจด้วย truth table -- ห้ามสรุปผลดังกล่าว
  \end{block}
\end{frame}


%% =====================================================================
\section{กรณีศึกษาและการพิสูจน์ด้วยตาราง}
%% =====================================================================

\begin{frame}{ตัวอย่าง: ระบบการตัดสินใจอนุมัติสินเชื่อ}
  \begin{block}{เงื่อนไข}
    \footnotesize
    1) รายได้ $>$ 50,000 และไม่มีหนี้ค้าง $\Rightarrow$ อนุมัติ \quad
    2) รายได้ $>$ 50,000 แต่มีหนี้ค้าง $\Rightarrow$ พิจารณาเพิ่ม \quad
    3) รายได้ $\leq$ 50,000 $\Rightarrow$ ไม่อนุมัติ
  \end{block}
  ให้ $I=$ รายได้สูง, \ $D=$ มีหนี้ค้าง, \ $A=$ ได้รับอนุมัติ \quad $\Rightarrow$ \quad $A \equiv I \land \neg D$
  \begin{center}\small
  \begin{tabular}{|c|c|c|c|}
    \hline
    $I$ & $D$ & $\neg D$ & $A \equiv I \land \neg D$ \\ \hline
    \TT & \TT & \FF & \FF \\ \hline
    \TT & \FF & \TT & \TT \\ \hline
    \FF & \TT & \FF & \FF \\ \hline
    \FF & \FF & \TT & \FF \\ \hline
  \end{tabular}
  \end{center}
  \footnotesize อนุมัติเฉพาะกรณีรายได้สูงและไม่มีหนี้ค้างเท่านั้น ($I=$T, $D=$F)
\end{frame}

\begin{frame}{ทฤษฎีบท: $p \rightarrow q \equiv \neg p \lor q$}
  พิสูจน์ด้วยตารางค่าความจริง
  \begin{center}\small
  \begin{tabular}{|c|c|c|c|c|c|}
    \hline
    \textbf{$p$} & \textbf{$q$} & \textbf{$p\rightarrow q$} & \textbf{$\neg p$} & \textbf{$\neg p \lor q$} & \textbf{การสมมูลกัน} \\ \hline
    \TT & \TT & \TT & \FF & \TT & \checkmark \\ \hline
    \TT & \FF & \FF & \FF & \FF & \checkmark \\ \hline
    \FF & \TT & \TT & \TT & \TT & \checkmark \\ \hline
    \FF & \FF & \TT & \TT & \TT & \checkmark \\ \hline
  \end{tabular}
  \end{center}
  \vspace{0.5em}
  คอลัมน์ $p\rightarrow q$ และ $\neg p\lor q$ เหมือนกันทุกแถว $\Rightarrow$ สมมูลกัน
  \begin{exampleblock}{ความสำคัญ}
    ใช้แปลงประพจน์เงื่อนไขให้อยู่ในรูปที่ใช้เฉพาะ AND, OR, NOT
    ซึ่งจำเป็นต่อการออกแบบวงจรดิจิทัลและการเขียนโปรแกรม
  \end{exampleblock}
\end{frame}

\begin{frame}{การใช้ตารางค่าความจริงในการพิสูจน์}
  ตารางค่าความจริงใช้ตรวจสอบได้ 3 วัตถุประสงค์หลัก
  \begin{enumerate}
    \item ประพจน์เป็น \textbf{สัจนิรันดร์} หรือไม่ -- ดูว่าคอลัมน์ผลลัพธ์เป็น \TT\ ทุกแถวหรือไม่
    \item ประพจน์สองประพจน์ \textbf{สมมูลกัน} หรือไม่ -- เปรียบเทียบคอลัมน์ผลลัพธ์
    \item การให้เหตุผล \textbf{สมเหตุสมผล} หรือไม่ -- เขียน (เหตุทั้งหมด $\land$) $\rightarrow$ ผล
          แล้วตรวจว่าเป็นสัจนิรันดร์หรือไม่
  \end{enumerate}
  \vspace{0.6em}
  \begin{block}{ข้อดี--ข้อจำกัด}
    ให้ผลแน่นอนและตรวจสอบได้เสมอ แต่จำนวนแถว $2^n$ โตเร็วเมื่อตัวแปรเพิ่มขึ้น
  \end{block}
\end{frame}


%% =====================================================================
\section{กิจกรรมในชั้นเรียน (แบบฝึกหัดทำในห้อง)}
%% =====================================================================

\begin{frame}{วิธีทำกิจกรรม}
  \begin{block}{คำชี้แจง}
    ให้นักศึกษาทำกิจกรรมต่อไปนี้ \textbf{ในชั้นเรียน} ลงในสมุด/กระดาษ
    ภายในเวลาที่กำหนด แล้วร่วมเฉลยและอภิปรายพร้อมกัน
  \end{block}
  \vspace{0.5em}
  \begin{itemize}
    \item ทำเดี่ยวหรือจับคู่ก็ได้ (ตามที่อาจารย์กำหนด)
    \item เมื่อหมดเวลา อาจารย์จะคลิกเพื่อ \textbf{เฉลย} ทีละข้อ
    \item อาสาสมัครออกมาอธิบายวิธีคิดหน้าชั้น
  \end{itemize}
  \vspace{0.5em}
  \begin{exampleblock}{เกณฑ์}
    เน้นกระบวนการคิดและการให้เหตุผล ไม่ใช่แค่คำตอบสุดท้าย
  \end{exampleblock}
\end{frame}

\begin{frame}{กิจกรรมที่ 1: การอนุมาน \eng{6 นาที}}
  จงระบุว่าใช้กฎอนุมานใด และหาผลสรุป
  \begin{enumerate}
    \item ถ้าวันนี้ฝนตก แล้วฉันจะไม่ออกไปข้างนอก ; วันนี้ฝนตก \quad $\dots$ ?
    \item ถ้าข้าพเจ้าสอบผ่าน แล้วข้าพเจ้าจะดีใจ ; ข้าพเจ้าไม่ดีใจ \quad $\dots$ ?
    \item ถ้า A แล้ว B ; ถ้า B แล้ว C \quad $\dots$ ?
    \item ระบบจะแจ้งเตือนทางอีเมลหรือ SMS ; ไม่ได้แจ้งทางอีเมล \quad $\dots$ ?
  \end{enumerate}
  \pause
  \begin{exampleblock}{เฉลย}
    \footnotesize
    1) Modus Ponens $\rightarrow$ ``ฉันจะไม่ออกไปข้างนอก'' \quad
    2) Modus Tollens $\rightarrow$ ``ข้าพเจ้าสอบไม่ผ่าน'' \\
    3) Hypothetical Syllogism $\rightarrow$ ``ถ้า A แล้ว C'' \quad
    4) Disjunctive Syllogism $\rightarrow$ ``แจ้งทาง SMS''
  \end{exampleblock}
\end{frame}

\begin{frame}{กิจกรรมที่ 2: ตรวจความสมเหตุสมผล \eng{8 นาที}}
  \begin{block}{โจทย์}
    จงตรวจว่าการให้เหตุผลต่อไปนี้สมเหตุสมผลหรือไม่ ด้วยตารางค่าความจริง \\
    ``จากสมมติฐาน $p \rightarrow q$ และ $q$ สรุปว่า $p$''
  \end{block}
  \pause
  \begin{center}\scriptsize
  \begin{tabular}{|c|c|c|c|c|}
    \hline
    $p$ & $q$ & $p\rightarrow q$ & $(p\rightarrow q)\land q$ & $[(p\rightarrow q)\land q]\rightarrow p$ \\ \hline
    \TT & \TT & \TT & \TT & \TT \\ \hline
    \TT & \FF & \FF & \FF & \TT \\ \hline
    \rowcolor{lpink}\FF & \TT & \TT & \TT & \FF \\ \hline
    \FF & \FF & \TT & \FF & \TT \\ \hline
  \end{tabular}
  \end{center}
  \begin{alertblock}{สรุป}
    มีแถวที่ผลลัพธ์เป็น \FF\ (แถวที่ 3) $\Rightarrow$ \textbf{ไม่เป็นสัจนิรันดร์ จึงไม่สมเหตุสมผล} \\
    นี่คือข้อผิดพลาด \emph{Affirming the Consequent} (การยอมรับผล)
  \end{alertblock}
\end{frame}


%% ====================== สรุป + แบบฝึกหัด ======================
\begin{frame}{สรุปสัปดาห์ที่ 3}
  \begin{exampleblock}{1. ตัวเชื่อมเพิ่มเติม}
    XOR ($\oplus$), NAND ($\uparrow$), NOR ($\downarrow$) -- ใช้ NAND อย่างเดียวสร้างทุกตัวเชื่อมได้ (Universal Gate)
  \end{exampleblock}
  \begin{exampleblock}{2. ประพจน์เปิดและตัวบ่งปริมาณ}
    $\forall x\,P(x)$ (ทุกตัว), $\exists x\,P(x)$ (มีบางตัว) -- แสดง $\forall$ เท็จ ใช้ \textbf{ตัวอย่างค้าน} 1 ตัว
  \end{exampleblock}
  \begin{exampleblock}{3. กฎการอนุมาน (Rules of Inference)}
    Modus Ponens, Modus Tollens, Hypothetical Syllogism, Disjunctive Syllogism \\
    -- \textbf{พิสูจน์ด้วยตารางค่าความจริงว่าเป็นสัจนิรันดร์}
  \end{exampleblock}
  \begin{exampleblock}{4. Fallacy (ข้อผิดพลาด)}
    \textbf{อย่า} สรุปแบบ Affirming the Consequent / Denying the Antecedent (ตรวจด้วยตารางก็ไม่เป็นสัจนิรันดร์)
  \end{exampleblock}
\end{frame}

\begin{frame}{สรุปบทที่ 1 (ภาพรวมตรรกศาสตร์)}
  \begin{itemize}
    \item \textbf{ประพจน์} + \textbf{ตัวเชื่อม} $\neg,\land,\lor,\rightarrow,\leftrightarrow$ (+ XOR, NAND, NOR)
    \item \textbf{ตารางค่าความจริง} = เครื่องมือวิเคราะห์ขั้นพื้นฐาน ($2^n$ กรณี)
    \item \textbf{ประเภท}: สัจนิรันดร์ (จริงเสมอ), ข้อขัดแย้ง (เท็จเสมอ), ประพจน์กลาง
    \item \textbf{ความสมมูล} + กฎพีชคณิตของประพจน์ (เดอมอร์แกน ฯลฯ)
    \item \textbf{ตัวบ่งปริมาณ} $\forall,\exists$ และนิเสธของตัวบ่งปริมาณ
    \item \textbf{การอนุมาน}: Modus Ponens, Modus Tollens, Hypothetical \& Disjunctive Syllogism
    \item \textbf{การประยุกต์}: เงื่อนไขในโปรแกรมคอมพิวเตอร์, วงจรอิเล็กทรอนิกส์ (Logic Gate)
  \end{itemize}
  \begin{exampleblock}{}
    \centering ตรรกศาสตร์เป็น \textbf{รากฐานสำคัญ} ของการออกแบบระบบคอมพิวเตอร์และระบบตัดสินใจ
  \end{exampleblock}
\end{frame}

\begin{frame}{แบบฝึกหัดมอบหมาย -- สัปดาห์ที่ 3 (ส่งสัปดาห์หน้า)}
  \begin{block}{คำชี้แจง}
    ทำลงใน \textbf{กระดาษ A4} แสดงวิธีทำให้ชัดเจน -- รวม 5 ข้อ คะแนนรวม \textbf{10 คะแนน}
  \end{block}
  \vspace{0.3em}
  \begin{enumerate}\small
    \item \textbf{[2 คะแนน]} จงพิสูจน์ว่า Hypothetical Syllogism เป็นสัจนิรันดร์
          ด้วยตารางค่าความจริง: $[(p\rightarrow q)\land(q\rightarrow r)]\rightarrow(p\rightarrow r)$
    \item \textbf{[2 คะแนน]} จงพิสูจน์ว่า Disjunctive Syllogism เป็นสัจนิรันดร์
          ด้วยตารางค่าความจริง: $[(p\lor q)\land \neg p]\rightarrow q$
    \item \textbf{[2 คะแนน]} จงเขียน $p \lor q$ และ $p \rightarrow q$ โดยใช้ NAND ($\uparrow$) อย่างเดียว
          พร้อมแสดงการพิสูจน์ด้วยตารางค่าความจริงว่าสมมูลกัน
    \item \textbf{[2 คะแนน]} กำหนดเอกภพ $U=\{1,2,3,4,5\}$ และประพจน์ $P(x)$: ``$x$ หารด้วย 2 ลงตัว''
          \begin{itemize}\small
            \item จงหาค่าความจริงของ $\forall x\,P(x)$ และ $\exists x\,P(x)$
            \item จงเขียนนิเสธของประพจน์ทั้งสอง
            \item ถ้า $\forall x\,P(x)$ เป็นเท็จ ให้ยกตัวอย่างค้าน
          \end{itemize}
    \item \textbf{[2 คะแนน -- บังคับ]} ระบบรับสมัครงาน: ผู้สมัครผ่านการคัดเลือกถ้า
          \textbf{(อายุ $\geq$ 22 และ GPA $\geq$ 3.0) หรือ (ประสบการณ์ $\geq$ 2 ปี)}
          \begin{itemize}\small
            \item กำหนดสัญลักษณ์: $A=$ อายุ$\geq$22, $G=$ GPA$\geq$3.0, $E=$ ประสบการณ์$\geq$2
            \item เขียนเงื่อนไขการผ่านคัดเลือกเป็นสูตรตรรกศาสตร์
            \item สร้างตารางค่าความจริงของเงื่อนไขดังกล่าว
          \end{itemize}
  \end{enumerate}
\end{frame}

\begin{frame}[plain]
  \vfill
  \centering
  {\Large\color{navy}\textbf{จบบทที่ 1 : ตรรกศาสตร์ (Logic)}}\par
  \vspace{0.8em}
  {\color{gray} สัปดาห์ถัดไป: บทที่ 2 : เซต (Set Theory)}
  \vfill
\end{frame}
